\documentclass[12pt]{article}
\usepackage{fontspec}
\usepackage{polyglossia}
\setmainlanguage{russian}
\defaultfontfeatures{Ligatures=TeX}
\setmainfont{Times New Roman}
\usepackage{geometry}
\geometry{
	a4paper,
	left=25mm,
	right=20mm,
	top=10mm,
	bottom=20mm
}

\usepackage{setspace}
\onehalfspacing
\usepackage{parskip}
\usepackage{microtype}
\usepackage{ragged2e}
\usepackage{multicol}
\usepackage{fancyhdr}

% ---------- Колонтитулы ----------
\pagestyle{fancy}
\fancyhf{}
\renewcommand{\headrulewidth}{0pt}
\fancyfoot[C]{\thepage}

%-------------- запрет переносов -------
\sloppy      
\hyphenpenalty=10000
\exhyphenpenalty=100
\pretolerance=1000
\tolerance=2000

% ---------- Макросы ----------
% Отправитель (фиксированный блок)
\newcommand{\Sender}{%
\noindent\textbf{Индивидуальный предприниматель} \\
\textbf{Мраморнов Александр Вячеславович} \\
%\linebreak
г. Краснодар, с/т № 2 А/О «Югтекс» \\
ул. Зеленая, 472 \\
ИНН 231200665168 \\
ОГРНИП 310231220400043\\
Телефон: 8-918-451-66-11 \\
e-mail: 4516611@gmail.com \\
\linebreak
Исх. № б/н \,\, \LetterDate
}

% Получатель (шаблон)
\newcommand{\Recipient}{%
		\textbf{}\\
	\textbf{} \\
	\\
\\
 \\
	\linebreak
         электронно: borodkinai@mail.ru\\
         8-910-985-55-11
}

% Переменные
\newcommand{\LetterDate}{\today}

%%%%%%%%%%!!!
\newcommand{\ReferenceDate}{21.04.2026}
%%%%%%%%%%!!!
\newcommand{\Предмет}{ }

\begin{document}
	
	% ВЕРХНИЙ БЛОК: слева отправитель, справа получатель
	\noindent
	\begin{minipage}[t]{75mm}
		\Sender
	\end{minipage}%
	\hfill
	\begin{minipage}[t]{75mm}
		\Recipient
	\end{minipage}
	
	\vspace{10mm}
	
	% Заголовок
	\begin{center}
		\Large \textbf{Информационное письмо}
	\end{center}
	
	\vspace{2mm}
	
На Ваш запрос  от \ReferenceDate\ сообщаю о возможности проведения судебной автотехнической экспертизы.
В объеме поставленных вопросов:\\
1. 
Имеются ли неисправности блока цилиндров
WUAZZZ8T5EA900137, номер 079103023S?
В случае положительного ответа на предыдущий
неисправностей?
Подлежит ли ремонту (восстаноњлению) блок цилиндра?
транспортного
вопрос, указать
средства
причины возникновения
В случае положительного ответа на предыдущий вопрос рассчитать стоимость восстановительного
ремонта.


 стоимость работ составит {10 000 (Десять тысяч) рублей}. \\
Производство исследования будет поручено специалисту:\\ 
{Мраморнову Александру Вячеславовичу}, имеющему высшее техническое образование по специальности «техническая физика», квалификация – инженер-физик, специальное образование в области независимой технической экспертизы транспортных средств: Диплом ПП-I № 424167, квалификация: эксперт-техник (специализация 150210 специальности 190601.65 – Автомобили и автомобильное хозяйство), состоящий в Государственном реестре экспертов-техников (№ в реестре 256), стаж экспертной работы 18 лет, общий трудовой стаж с 1989 г.
	
\vspace{8mm}
	
Индивидуальный предприниматель \\
Мраморнов А.В.  \hfill {\rule{6cm}{0.4pt}
	
\end{document}

